\section{Introduction}

A model that reasons about how words sound needs a written form for sound. The common choice is the International Phonetic Alphabet, which is accurate, widely used among linguists, and holds a large amount of recorded pronunciation data. For a model pipeline it raises a few practical problems.

A single sound is often a base letter followed by one or more combining diacritics, so the number of characters does not match the number of sounds. Some of the newer symbols sit in the supplementary planes of Unicode, which many tokenizers handle as several units. And the subword or byte tokenizers that most models use break one pronunciation into several tokens that have no fixed relation to the phones. A model then has to learn that a scattered set of fragments stands for one sound.

Talk is a phonetic encoding designed around these problems. It has two written forms. The first is a plain ASCII spelling, where each sound is a base consonant or vowel followed by modifiers, so a person can type pronunciations on a normal keyboard and a program can read them without special handling. The second is a machine form, where each canonical sound is a single code point, so a word becomes a compact string of one token per sound. A short example, in Python:

\begin{lstlisting}
import talk

# prajna, written in IPA
talk.talk_to_ipa("pradjy~ny~a_") # => "pradʑɲaː"

# one Hangul code point per sound
talk.machine("pradjy~ny~a_") # => a 7 code point string

# the sounds, each with its features
talk.segment("pradjy~ny~a_") # => [p, r, a, d, jy~, ny~, a_]

# and the reverse, IPA back to talk
talk.ipa_to_talk("prɐdʑɲaː") # => "pradjy~ny~a_"
\end{lstlisting}

The contribution is a small set of design choices that work together.

\begin{itemize}
\item A phonetic alphabet in plain ASCII, where each sound is a base plus modifiers in a fixed order, so every sound has exactly one spelling.
\item A machine form that maps each canonical sound to one code point, so a word is one token per sound. The Hangul syllable block holds the individual sounds, and Chinese characters hold consonant and vowel clusters.
\item A fixed, append-only map from sound to code point, so a token never changes meaning across releases.
\item Conversion to and from IPA, without loss in the talk direction, and splitting of a word into syllables.
\item Implementations for Python and TypeScript with no runtime dependencies.
\end{itemize}

Talk sits alongside earlier work. X-SAMPA gives an ASCII spelling of IPA for keyboard entry \cite{wells1995sampa}, and tools such as Epitran map orthography to IPA for many languages \cite{mortensen2018epitran}. Subword methods such as byte-pair encoding and unigram models are how model vocabularies are usually built \cite{sennrich2016bpe, kudo2018sentencepiece}. Talk aims instead at a phone-level and cluster-level vocabulary directly, with a fixed one-token-per-sound map and a bridge to IPA in both directions.

The intended readers build text-to-speech, speech recognition, grapheme-to-phoneme, or multilingual systems, and want a compact and stable phonetic vocabulary in place of raw IPA. The source is available \cite{cluesurf2026talk}.
